\chapter{Keliling dan Luas}\label{ch:g5:areas}

Kelas 4 membilang petak (\cref{ch:g4:measure}); bab ini memperoleh
\emph{rumus} luas yang pertama --- panjang kali lebar untuk persegi
panjang --- dan mempelajari satuan cm$^2$ dan m$^2$. Rumus untuk
segitiga dan bangun lain matang pada \cref{ch:g6:measure} dan
\cref{ch:g7:areas}.

\section{Dua ukuran yang berbeda}

\begin{example}[Tepi lawan permukaan]\label{ex:g5:areas:contrast}
Seorang tukang kebun memerlukan \emph{\omterm{def:g4:measure:perimeter}{keliling}} untuk membeli pagar
dan \emph{luas} untuk membeli benih rumput. Keduanya tidak saling
mengikuti: menarik persegi panjang menjadi lebih tipis dan lebih
panjang dapat mempertahankan luasnya sementara \omterm{def:g4:measure:perimeter}{kelilingnya} bertambah
--- bandingkan persegi panjang $6 \times 2$ dan $4 \times 3$ (luasnya
sama, $12$, \omterm{def:g4:measure:perimeter}{kelilingnya} $16$ dan $14$).
\end{example}

\section{Rumus persegi panjang}

\begin{proposition}[Luas persegi panjang]\label{prop:g5:areas:rectangle}
Persegi panjang dengan panjang $L$ dan lebar $w$ (dalam satuan yang sama) punya luas
\[
A = L \times w ,
\]
dalam satuan persegi: cm$^2$ (persegi dengan sisi $1$ cm), m$^2$
(persegi dengan sisi $1$ m), \dots
\end{proposition}

\begin{proof}[Mengapa, dengan membilang]
Tutupilah persegi panjang itu dengan persegi satuan: $w$ baris
berisi $L$ persegi, jadi seluruhnya $L \times w$ persegi --- persis
persegi panjang perkalian pada \cref{ch:g3:mult}.
\end{proof}

\begin{omfigure}
\begin{tikzpicture}[scale=0.62]
  \draw[gray!50, thin] (0,0) grid (6,3);
  \draw[very thick, omDef] (0,0) rectangle (6,3);
  \fill[omThm!30] (0,2) rectangle (1,3);
  \draw[omThm, thick] (0,2) rectangle (1,3);
  \node[right, font=\small, omThm] at (1.15,2.5) {$1$ cm$^2$};
  \node[below, font=\small] at (3,-0.15) {$6$ cm};
  \node[left, font=\small] at (0,1.5) {$3$ cm};
  \node[font=\small] at (3.9,1.2) {$6 \times 3 = 18$ cm$^2$};
\end{tikzpicture}

{\small Tiga baris berisi enam persegi sentimeter: rumus
$L \times w$ adalah hitungan itu, ditulis sekali untuk selamanya.}
\end{omfigure}

\begin{example}\label{ex:g5:areas:rectangle}
Sebuah permadani berukuran $2.5$ m kali $2$ m: luasnya
$2.5 \times 2 = 5$ m$^2$. Sebuah perangko berukuran $3$ cm kali
$2.4$ cm: luasnya $3 \times 2.4 = 7.2$ cm$^2$. Rumusnya sama untuk
satuan apa pun --- asalkan kedua sisinya memakai satuan yang
\emph{sama}.
\end{example}

\begin{example}[Satuan luas berubah dengan seratus]\label{ex:g5:areas:units}
$1$ m $= 100$ cm, tetapi $1$ m$^2 = 100 \times 100 = 10\,000$ cm$^2$:
satu meter persegi adalah petak $100$ kali $100$ sentimeter persegi.
Satuan luas melompat dengan ratusan, bukan dengan puluhan.
\end{example}

\section{Bangun gabungan}

\begin{method}[Potong, jumlahkan, kurangkan]\label{met:g5:areas:composite}
Untuk bangun yang tersusun dari persegi panjang (huruf L, huruf T, sebuah bingkai):
\begin{enumerate}
  \item potonglah menjadi beberapa persegi panjang, atau lengkapi
        menjadi satu persegi panjang yang besar;
  \item hitunglah luas tiap persegi panjang dengan rumusnya;
  \item jumlahkan potongannya --- atau kurangkan lubangnya dari
        persegi panjang yang besar tadi;
  \item periksalah dengan membilang petaknya secara kasar.
\end{enumerate}
\end{method}

\begin{example}\label{ex:g5:areas:composite}
Sebuah ruangan berbentuk L: persegi panjang $6$ m $\times$ $4$ m
dengan pojok $2$ m $\times$ $2$ m yang hilang.
\[
\text{Luas} = 6 \times 4 - 2 \times 2 = 24 - 4 = 20 \text{ m}^2 .
\]
Atau potonglah huruf L itu menjadi jalur $6 \times 2$ dan jalur
$4 \times 2$: $12 + 8 = 20$ m$^2$ --- dua jalan, satu jawaban.
\end{example}

\begin{example}[Setengah persegi panjang]\label{ex:g5:areas:halfrect}
Memotong sebuah persegi panjang sepanjang diagonalnya memberi dua
segitiga yang sama luasnya: masing-masing adalah \emph{setengah} persegi
panjang itu. Segitiga siku-siku dengan sisi siku-siku $6$ cm dan
$4$ cm karena itu berluas $\frac{6 \times 4}{2} = 12$ cm$^2$ ---
gambar yang layak diingat untuk \cref{ch:g6:measure}, tempat ia
menjadi sebuah rumus.
\end{example}

\begin{omfigure}
\begin{tikzpicture}[scale=0.62]
  \draw[gray!50, thin] (0,0) grid (6,4);
  \draw[very thick] (0,0) rectangle (6,4);
  \fill[omDef!25] (0,0) -- (6,0) -- (0,4) -- cycle;
  \draw[very thick, omDef] (6,0) -- (0,4);
  \node[font=\small] at (1.7,1.2) {setengah};
  \node[font=\small] at (4.3,2.8) {setengah};
\end{tikzpicture}

{\small Diagonalnya memotong persegi panjang $6 \times 4$ menjadi
dua segitiga yang sama besar, masing-masing $12$ petak.}
\end{omfigure}

\section{Latihan}

\begin{exercise}[$\star$]\label{exo:g5:areas:1}
Hitunglah luas dan \omterm{def:g4:measure:perimeter}{keliling} persegi panjang $8$ cm $\times$ $5$ cm.
Jawaban mana yang dalam cm, mana yang dalam cm$^2$?
\end{exercise}

\begin{exercise}[$\star$]\label{exo:g5:areas:2}
Hitunglah luasnya: persegi dengan sisi $9$ cm; persegi panjang
$12$ m $\times$ $7$ m; persegi panjang $4.5$ cm $\times$ $6$ cm.
\end{exercise}

\begin{exercise}[$\star$]\label{exo:g5:areas:3}
Sebuah persegi panjang berluas $63$ cm$^2$ dan panjangnya $9$ cm.
Carilah lebarnya, lalu \omterm{def:g4:measure:perimeter}{kelilingnya}.
\end{exercise}

\begin{exercise}[$\star$]\label{exo:g5:areas:4}
Gambarlah dua persegi panjang berbeda yang luasnya $24$ petak di
kertas berpetak, lalu hitunglah kedua \omterm{def:g4:measure:perimeter}{kelilingnya}. Luasnya sama --- \omterm{def:g4:measure:perimeter}{kelilingnya} sama juga?
\end{exercise}

\begin{exercise}[$\star$]\label{exo:g5:areas:5}
Ubahlah: $3$ m$^2$ menjadi cm$^2$; \quad $50\,000$ cm$^2$ menjadi
m$^2$. (Ingat \cref{ex:g5:areas:units}: dengan ratusan!)
\end{exercise}

\begin{exercise}[$\star$]\label{exo:g5:areas:6}
Sebuah bangun berbentuk huruf T tersusun dari batang mendatar
$8 \times 2$ di atas batang tegak $2 \times 5$ (dalam cm).
Hitunglah luasnya dengan menjumlahkan dua persegi panjang.
\end{exercise}

\begin{exercise}[$\star$]\label{exo:g5:areas:7}
Sebuah bingkai foto: persegi panjang $30$ cm $\times$ $20$ cm dengan
jendela persegi panjang $24$ cm $\times$ $14$ cm yang dilubangi.
Berapa luas kayu yang dipakai bingkai itu?
\end{exercise}

\begin{exercise}[$\star$]\label{exo:g5:areas:8}
Hitunglah luas segitiga siku-siku dengan sisi siku-siku $8$ cm dan
$6$ cm (setengah persegi panjang, \cref{ex:g5:areas:halfrect}).
\end{exercise}

\begin{exercise}[$\star$]\label{exo:g5:areas:9}
Sebuah petak sayur berbentuk persegi panjang berukuran $7$ m kali
$4$ m. Benihnya berharga $2$ per meter persegi, dan pagarnya $3$
per meter. Hitunglah biaya benihnya, lalu biaya pagarnya.
\end{exercise}

\begin{exercise}[$\star\star$]\label{exo:g5:areas:10}
Lantai sebuah lorong panjangnya $12$ m dan lebarnya $2$ m, dan harus
ditutup ubin persegi dengan sisi $50$ cm. Berapa ubin yang
diperlukan? (Ubahlah satuannya dulu, atau bilanglah ubinnya ke tiap arah.)
\end{exercise}

\begin{exercise}[$\star\star$]\label{exo:g5:areas:11}
Lipatduakan sisi persegi panjang $4$ cm $\times$ $3$ cm. Apa yang
terjadi pada \omterm{def:g4:measure:perimeter}{kelilingnya}? Pada luasnya? (Hitunglah keduanya sebelum
dan sesudah --- kedua jawabannya berbeda!)

\end{exercise}
